\section{After Measurements}
This section describes the considerations that should be made after measurements have been performed. At its core, this involves reflecting on the procedure, what has been measured, and how those measurements compare to what was intended and expected.

It is not uncommon to realize that the measurements may not represent what was originally intended, or that they are only weakly linked to the target or goal. This is also true for exploratory measurements, where reflection may reveal that alternative setups or different quantities would better describe the phenomenon under investigation.


\subsection{Data Storage and Accessibility}
Once all data described in the procedure has been collected, it is imperative that it is stored securely and remains accessible. If cloud services are used, it is good practice to have another person verify access to the data to ensure that it is available and has not been corrupted.

It is important to recognize that high-quality measurement data is highly valuable. For example, a reference laboratory may perform only a limited number of measurements on a 10 V DC standard and issue a signed, accredited certificate at a cost of several thousand euros. 

The collected data should therefore be treated as a valuable asset and stored accordingly, with appropriate backup and access measures in place.

\subsection{Analyze the Data}
The collected data should be carefully analyzed, and the presence of outliers should be evaluated. If outliers are identified, it is essential to determine whether they arise from measurement errors, data entry mistakes, or other identifiable causes.

If an outlier is to be excluded, a clear and well-justified reason must be provided. Ideally, the removal of the outlier should be documented, along with the rationale for its exclusion.

It is important to recognize that the data has been observed under the defined measurement procedure. If the procedure is valid, the origin of an outlier may not be immediately clear-it could be a genuine feature of the phenomenon under investigation, or it could result from misuse of the measurement setup, data entry errors, or other issues.

Data must not be excluded simply because it appears inconsistent or does not fit expectations. Such practice is not scientifically valid, may lead to incorrect conclusions, and is both unethical and in bad faith.

\subsection{General Estimation of Uncertainty}
Depending on the requirements and agreements made with the supervisor, a general 95\% coverage uncertainty (k = 1.96) should be estimated and included in the dataset. It is recognized that, for some scientific work, uncertainty estimation may not always be the primary focus. However, at the core of science lies metrology, and any well-prepared dataset should include some degree of uncertainty estimation.

In many cases, a simplified approach may be sufficient. For example, if measurements are performed using a vector network analyzer under consistent conditions-such as similar attenuator settings, mixer levels, reference levels, and bandwidth-it may be reasonable to assign a uniform uncertainty, e.g., $\pm$ 2 dB, to all measurements.

Such an estimate can be valid, provided it is justified. The basis for the estimation should always be clearly described, even if a simplified method is used. I would like to quote the international recognized bible (BIPM GUM) within uncertainties here:

\begin{quote}
\centering
\textbf{In general, the result of a measurement is only an approximation or estimate of the value of the measurand and is complete only when accompanied by a statement of the uncertainty of that estimate.}
\end{quote}

Please be aware that the evaluation of uncertainty is not something that should be delegated to technicians. While they may assist with instrument setup and configuration, the evaluation and estimation of uncertainty is the responsibility of the experimenter and the designer of the measurement procedure.

There are numerous documents and general guides available to support this process, several of which are referenced throughout this document.


\subsection{Review and Update Measurement Procedure}
The final step is to review the measurement procedure and evaluate whether any elements should be removed, added, or modified. The procedure must always describe the exact method by which the data was collected.

If the procedure is changed and the collected data no longer aligns with the documented method, the data is no longer valid under that procedure. In such cases, the data must either be collected again or documented under a separate procedure specific to that dataset.

There is no wiggle room in this principle: data must always follow the documented method. This is the cornerstone of reliable and traceable measurements, ensuring that every dataset can be understood, reproduced, and validated.

It must be made absolutely clear: under no circumstances should data or procedures be altered to make results appear more favorable. Any attempt to manipulate measurements undermines the scientific method and the reliability of the dataset. The integrity of the measurement process is paramount; it is not subject to personal preference, intuition, or aesthetic judgment. 

\subsection{Publication}
Once all the preceding steps have been completed and both the experimenter and supervisor are satisfied, the data is ready for publication. If this document has been followed thoroughly, the resulting dataset should be of such quality that it rivals that of the finest laboratories in industry. 

High-quality measurements are a source of pride, and producing them requires effort, diligence, and attention to detail. Following rigorous procedures ensures that the data is trustworthy, reproducible, and valuable to others in the scientific and engineering community.